% Generated draft: requires mathematical and source review.
\documentclass[11pt]{article}
\usepackage[margin=25mm]{geometry}
\usepackage{fontspec}
\setmainfont{Latin Modern Roman}
\usepackage{amsmath,amssymb,mathtools,unicode-math}
\setmathfont{Latin Modern Math}
\usepackage{enumitem,needspace,xurl,hyperref}
\setlength{\parindent}{0pt}
\setlength{\parskip}{5pt}
\setlength{\emergencystretch}{3em}
\newcommand{\N}{\mathbb{N}}
\newcommand{\Z}{\mathbb{Z}}
\newcommand{\Q}{\mathbb{Q}}
\newcommand{\R}{\mathbb{R}}
\newcommand{\C}{\mathbb{C}}
\newcommand{\F}{\mathbb{F}}
\newcommand{\A}{\mathbb{A}}
\newcommand{\PP}{\mathbb{P}}
\newcommand{\E}{\mathbb{E}}
\newcommand{\eps}{\varepsilon}
\newcommand{\dd}{\,\mathrm{d}}
\begin{document}
\section*{UnsolvedMath ID: 30000435}
\textbf{Source problem code:} catalogue entry
\subsection{Definitions and mathematical statement}
In the usual three-dimensional unfolding setting, let \(P\subset\R^3\) be a convex polytope and equip its boundary \(S=\partial P\) with its intrinsic piecewise-Euclidean metric. Let \(T\) be the intended intrinsic Delaunay triangulation, using the polytope vertices and geodesic edges. The usual local Delaunay condition is that the angles opposite an interior triangulation edge sum to at most \(\pi\). The source does not specify the dimension, vertex set, or intrinsic versus extrinsic convention explicitly; the display describes the standard surface interpretation.

The question is whether one can choose cuts contained in the one-skeleton \(T^{(1)}\) so that the connected cut-open surface develops into the Euclidean plane without overlap. A standard precise target is a spanning cut tree \(C\subseteq T^{(1)}\) through all polytope vertices, for which the metric completion \(S_C\) of the cut-open boundary is a disk and has a facewise isometric development
\[
\operatorname{dev}:S_C\longrightarrow\R^2
\]
with disjoint interiors of distinct developed triangles. The cut set is a selected subset of triangulation edges, not all edges, and the result is one connected net. Whether boundary contacts are allowed or full injectivity including the boundary is required must be confirmed; the usual non-overlapping-net convention permits boundary contacts while forbidding interior overlap.

The particular source question asks whether, when all faces of \(P\) are acute triangles, such a non-overlapping unfolding exists using only original polytope edges. Acute triangular faces satisfy the usual local intrinsic Delaunay angle condition, so this is the stated edge-cut special case. No cuts through face interiors, separate disconnected pieces, or general unrestricted surface-cut unfolding are substituted for that requirement.

For a nonunique Delaunay triangulation, the source does not state whether a favorable choice of \(T\) is allowed or every choice must admit the desired cuts. The spanning-tree, connected-disk interpretation also requires matching to the report's exact notion of unfolding. The supplied assessment checked 2026-08-21 reports no general non-overlapping theorem for this Delaunay-restricted or acute-face edge-restricted question.
\subsection{Short English statement}
Can a convex polyhedral surface be cut along selected intrinsic Delaunay edges into one planar net without interior overlap? The special case requires only original edges when all faces are acute triangles; cut and triangulation conventions need confirmation.
\subsection{Sources}
\begin{itemize}
\item Local UnsolvedMath registry, permanent ID 30000435, catalogue entry, \emph{Delaunay Unfoldings of Convex Polytope Boundaries}.
\item \emph{Discrete Differential Geometry} (2006), Oberwolfach Reports, pp. 653--728, DOI \url{https://doi.org/10.4171/owr/2006/12}.
\item Supplied reference labeled ``convex polytope unfolding status'': \url{https://jeffe.web.engr.illinois.edu/open/unfold.html}.
\item Supplied source-verification claim and literature triage checked 2026-08-21; no general unfolding theorem or full convention list is reproduced.
\end{itemize}
\end{document}
